\subsection{Tensor Product Bases and Maps}

\begin{theorem} \label{thm:tensor-bases-dimension-and-rank} \hlt{Tensor Bases, Dimension, and Coordinate Rank}\\
Let $\mathcal{U}$ and $\mathcal{V}$ be finite-dimensional with bases $\beta=(\mathbf{u}_1,\ldots,\mathbf{u}_m)$ and $\gamma=(\mathbf{v}_1,\ldots,\mathbf{v}_n)$.
\begin{enumerate}[label=(\roman*)]
\setlength{\itemsep}{0pt}
\item The tensors $\mathbf{u}_i\otimes\mathbf{v}_j$ form a basis of $\mathcal{U}\otimes\mathcal{V}$.
\item $\dim(\mathcal{U}\otimes\mathcal{V})=\dim(\mathcal{U})\dim(\mathcal{V})$.
\item Every $\mathbf{z}\in\mathcal{U}\otimes\mathcal{V}$ has a unique matrix $R_{\beta,\gamma}(\mathbf{z})=[r_{ij}]\in\operatorname{Mat}(m,n,\mathbb{F})$ such that $\mathbf{z}=\sum\limits_{i=1}^{m}\sum\limits_{j=1}^{n} r_{ij}(\mathbf{u}_i\otimes\mathbf{v}_j)$.\\
If $\beta',\gamma'$ are other bases, then $\operatorname{rank}R_{\beta',\gamma'}(\mathbf{z})=\operatorname{rank}R_{\beta,\gamma}(\mathbf{z})$.
\end{enumerate}
\begin{proof}
\begin{enumerate}[label=(\roman*)]
\setlength{\itemsep}{0pt}
\item The coordinate tensor model of Definition \ref{def:finite-coordinate-tensor-product-model} has the universal property in Definition \ref{def:tensor-product-universal-property}.\\
Hence Corollary \ref{cor:tensor-product-unique-up-to-isomorphism} identifies $\mathcal{U}\otimes\mathcal{V}\cong\operatorname{Mat}(m,n,\mathbb{F})$ with $\mathbf{u}_i\otimes\mathbf{v}_j\mapsto E_{i,j}$, and $\{E_{i,j}\}$ is a basis.
\item Counting the $mn$ basis tensors gives $\dim(\mathcal{U}\otimes\mathcal{V})=mn=\dim(\mathcal{U})\dim(\mathcal{V})$.
\item The basis in (i) gives unique coefficients $r_{ij}$, hence $R_{\beta,\gamma}(\mathbf{z})$. If $\beta',\gamma'$ are other bases,\\
then $R_{\beta',\gamma'}(\mathbf{z})=P R_{\beta,\gamma}(\mathbf{z})Q$ for invertible change-of-basis matrices $P,Q$, so $\operatorname{rank}R_{\beta',\gamma'}(\mathbf{z})=\operatorname{rank}R_{\beta,\gamma}(\mathbf{z})$.
\end{enumerate}
\hlt{Structural Commentary}: Tensor rank is coordinate-free matrix rank.\\
\hlt{Key Insight}: Changing bases performs row and column changes, which preserve matrix rank.
\end{proof}
\end{theorem}

\begin{definition} \label{def:tensor-coordinate-matrix-and-rank} \hlt{Tensor Coordinate Matrix and Rank}\\
Let $\mathcal{U},\mathcal{V}$ be finite-dimensional, with bases $\beta=(\mathbf{u}_1,\ldots,\mathbf{u}_m)$ and $\gamma=(\mathbf{v}_1,\ldots,\mathbf{v}_n)$.
\begin{enumerate}[label=(\roman*)]
\setlength{\itemsep}{0pt}
\item For $\mathbf{z}=\sum\limits_{i=1}^{m}\sum\limits_{j=1}^{n} r_{ij}(\mathbf{u}_i\otimes\mathbf{v}_j)$, define the \hlt{coordinate matrix} $R_{\beta,\gamma}(\mathbf{z})\defeq[r_{ij}]$.
\item The \hlt{rank of the tensor} $\mathbf{z}$ is $\operatorname{rank}(\mathbf{z})\defeq\operatorname{rank}R_{\beta,\gamma}(\mathbf{z})$, independent of $\beta,\gamma$ by Theorem \ref{thm:tensor-bases-dimension-and-rank}.
\end{enumerate}
\end{definition}

\begin{theorem} \label{thm:tensor-zero-test-and-decomposable-rank} \hlt{Zero Tests and Decomposable Rank}\\
Let $\mathcal{U}$ and $\mathcal{V}$ be vector spaces.
\begin{enumerate}[label=(\roman*)]
\setlength{\itemsep}{0pt}
\item $\forall \mathbf{u}\in\mathcal{U},\mathbf{v}\in\mathcal{V}$, $\mathbf{0}\otimes\mathbf{v}=\mathbf{u}\otimes\mathbf{0}=\mathbf{0}$.
\item If $\mathbf{u}_1,\ldots,\mathbf{u}_r$ are linearly independent and $\sum\limits_{i=1}^{r}\mathbf{u}_i\otimes\mathbf{v}_i=\mathbf{0}$, then $\mathbf{v}_1=\cdots=\mathbf{v}_r=\mathbf{0}$.
\item If $\dim\mathcal{U},\dim\mathcal{V}<\infty$, then $\mathbf{z}\neq\mathbf{0}$ is decomposable $\Longleftrightarrow \operatorname{rank}(\mathbf{z})=1$.\\
More generally, $\operatorname{rank}(\mathbf{z})=\min\left\{r\ge0:\mathbf{z}=\sum\limits_{i=1}^{r}\mathbf{x}_i\otimes\mathbf{y}_i,\ \mathbf{x}_i\in\mathcal{U},\mathbf{y}_i\in\mathcal{V}\right\}$.
\end{enumerate}
\begin{proof}
\begin{enumerate}[label=(\roman*)]
\setlength{\itemsep}{0pt}
\item The tensor map $t:\mathcal{U}\times\mathcal{V}\to\mathcal{U}\otimes\mathcal{V}$ is bilinear, hence linear in each coordinate; therefore $t(\mathbf{0},\mathbf{v})=t(\mathbf{u},\mathbf{0})=\mathbf{0}$.
\item For each $k$, extend the dual basis functional on $\operatorname{span}(\mathbf{u}_1,\ldots,\mathbf{u}_r)$ to $\varphi_k\in\mathcal{U}'$ with $\varphi_k(\mathbf{u}_i)=\delta_{ki}$.\\
The bilinear map $(\mathbf{u},\mathbf{v})\mapsto\varphi_k(\mathbf{u})\mathbf{v}$ induces $T_k\in\LL(\mathcal{U}\otimes\mathcal{V},\mathcal{V})$.\\
Applying $T_k$ gives $\mathbf{0}=T_k\left(\sum_{i=1}^{r}\mathbf{u}_i\otimes\mathbf{v}_i\right)=\sum\limits_{i=1}^{r}\delta_{ki}\mathbf{v}_i=\mathbf{v}_k$.
\item Choose bases $\beta,\gamma$. If $\mathbf{z}=\mathbf{x}\otimes\mathbf{y}\neq\mathbf{0}$, then $R_{\beta,\gamma}(\mathbf{z})=[\mathbf{x}]_\beta[\mathbf{y}]_\gamma^{\mathsf{T}}$, so $\operatorname{rank}(\mathbf{z})=1$.\\
Conversely, if $\operatorname{rank}R_{\beta,\gamma}(\mathbf{z})=r$, write $R_{\beta,\gamma}(\mathbf{z})=\sum\limits_{i=1}^{r}a_i b_i^{\mathsf{T}}$ with columns $a_i,b_i$.\\
Translating $a_i,b_i$ back to vectors $\mathbf{x}_i,\mathbf{y}_i$ gives $\mathbf{z}=\sum\limits_{i=1}^{r}\mathbf{x}_i\otimes\mathbf{y}_i$ as a sum of $r$ decomposable tensors.\\
Conversely, any sum of $s$ decomposable tensors gives a sum of $s$ rank-one coordinate matrices, so $\operatorname{rank}(\mathbf{z})\le s$; hence the displayed minimum equals $\operatorname{rank}(\mathbf{z})$.
\end{enumerate}
\hlt{Structural Commentary}: Linear functionals can test one tensor coordinate at a time.\\
\hlt{Key Insight}: Tensor rank is matrix rank after choosing bases, but the number itself survives changing bases.
\end{proof}
\end{theorem}

\begin{theorem} \label{thm:tensor-linearizes-bilinear-maps} \hlt{Bilinearity Equals Linearity on a Tensor Product}\\
For $\mathbb{F}$-vector spaces $\mathcal{U},\mathcal{V},\mathcal{W}$, the map $\Phi:\operatorname{Bil}_{\mathbb{F}}(\mathcal{U},\mathcal{V};\mathcal{W})\to\LL(\mathcal{U}\otimes\mathcal{V},\mathcal{W})$, $\Phi(f)=\widetilde{f}$,\\
where $\widetilde{f}(\mathbf{u}\otimes\mathbf{v})=f(\mathbf{u},\mathbf{v})$, is an isomorphism.
\begin{proof}
The universal property gives existence and uniqueness of $\widetilde{f}$.\\
Define $\Psi:\LL(\mathcal{U}\otimes\mathcal{V},\mathcal{W})\to\operatorname{Bil}_{\mathbb{F}}(\mathcal{U},\mathcal{V};\mathcal{W})$ by $\Psi(L)\defeq L\circ t$.\\
Then $\Psi(\Phi(f))=\widetilde{f}\circ t=f$, and $\Phi(\Psi(L))=L$ by uniqueness of the mediating map for $L\circ t$. Hence $\Phi$ is bijective.\\
For $a\in\mathbb{F}$ and $f,g$ bilinear, $\Phi(af+g)(\mathbf{u}\otimes\mathbf{v})=af(\mathbf{u},\mathbf{v})+g(\mathbf{u},\mathbf{v})=(a\Phi(f)+\Phi(g))(\mathbf{u}\otimes\mathbf{v})$, so $\Phi$ is linear.\\
\hlt{Structural Commentary}: Tensor products package bilinear maps as ordinary linear maps.\\
\hlt{Key Insight}: To define a linear map out of $\mathcal{U}\otimes\mathcal{V}$, define a bilinear rule on $\mathcal{U}\times\mathcal{V}$.
\end{proof}
\end{theorem}

\begin{method} \label{meth:defining-linear-maps-on-tensor-products} \hlt{Defining Maps on Tensor Products}\\
To define $L\in\LL(\mathcal{U}\otimes\mathcal{V},\mathcal{W})$:
\begin{enumerate}[label=(\roman*)]
\setlength{\itemsep}{0pt}
\item define a candidate rule $f:\mathcal{U}\times\mathcal{V}\to\mathcal{W}$ by $f(\mathbf{u},\mathbf{v})$;
\item prove $f\in\operatorname{Bil}_{\mathbb{F}}(\mathcal{U},\mathcal{V};\mathcal{W})$;
\item set $L\defeq\Phi(f)$ from Theorem \ref{thm:tensor-linearizes-bilinear-maps}, so $L(\mathbf{u}\otimes\mathbf{v})=f(\mathbf{u},\mathbf{v})$.
\end{enumerate}
Do not assign arbitrary values on decomposable tensors: they span $\mathcal{U}\otimes\mathcal{V}$ but are usually linearly dependent.
\end{method}

\begin{theorem} \label{thm:tensor-product-of-linear-maps} \hlt{Tensor Product of Linear Maps}\\
Let $A\in\LL(\mathcal{U},\mathcal{U}')$ and $B\in\LL(\mathcal{V},\mathcal{V}')$.\\
There exists a unique $A\otimes B:\mathcal{U}\otimes\mathcal{V}\to\mathcal{U}'\otimes\mathcal{V}'$ such that $(A\otimes B)(\mathbf{u}\otimes\mathbf{v})=A\mathbf{u}\otimes B\mathbf{v}$.\\
Moreover, $(A_2A_1)\otimes(B_2B_1)=(A_2\otimes B_2)(A_1\otimes B_1)$, $I_{\mathcal{U}}\otimes I_{\mathcal{V}}=I_{\mathcal{U}\otimes\mathcal{V}}$.\\
In finite dimensions, $[A\otimes B]_{\beta\otimes\gamma}^{\beta'\otimes\gamma'}=[A]_{\beta}^{\beta'}\otimes[B]_{\gamma}^{\gamma'}$ for lexicographically ordered tensor-product bases, where the right-hand $\otimes$ is the matrix Kronecker product.
\begin{proof}
The map $(\mathbf{u},\mathbf{v})\mapsto A\mathbf{u}\otimes B\mathbf{v}$ is bilinear, so Theorem \ref{thm:tensor-linearizes-bilinear-maps} gives the unique $A\otimes B$.\\
The two functorial identities hold on decomposable tensors $((A_2\otimes B_2)(A_1\otimes B_1))(\mathbf{u}\otimes\mathbf{v})=A_2A_1\mathbf{u}\otimes B_2B_1\mathbf{v}$ and $(I_{\mathcal{U}}\otimes I_{\mathcal{V}})(\mathbf{u}\otimes\mathbf{v})=\mathbf{u}\otimes\mathbf{v}$, and decomposable tensors span. If $A\mathbf{u}_j=\sum_i a_{ij}\mathbf{u}'_i$ and $B\mathbf{v}_\ell=\sum_k b_{k\ell}\mathbf{v}'_k$, then $(A\otimes B)(\mathbf{u}_j\otimes\mathbf{v}_{\ell}) =\sum_{i,k}a_{ij}b_{k\ell}(\mathbf{u}'_i\otimes\mathbf{v}'_k)$, which is the Kronecker block-entry rule.\\
\hlt{Structural Commentary}: Tensoring maps acts independently in the two tensor coordinates.\\
\hlt{Key Insight}: The mixed product property for Kronecker products is the coordinate shadow of functoriality.
\end{proof}
\end{theorem}

\begin{theorem} \label{thm:tensor-product-of-map-spaces} \hlt{Tensor Product of Spaces of Linear Maps}\\
Let $\mathcal{U},\mathcal{V},\mathcal{X},\mathcal{Y}$ be $\mathbb{F}$-vector spaces. The assignment $(A,B)\mapsto A\otimes B$ induces a natural injective map\\
$\Theta:\LL(\mathcal{U},\mathcal{X})\otimes\LL(\mathcal{V},\mathcal{Y}) \to \LL(\mathcal{U}\otimes\mathcal{V},\mathcal{X}\otimes\mathcal{Y})$ with $\Theta(A\otimes B)(\mathbf{u}\otimes\mathbf{v})=A\mathbf{u}\otimes B\mathbf{v}$.\\
If $\mathcal{U},\mathcal{V},\mathcal{X},\mathcal{Y}$ are finite-dimensional, then $\Theta$ is an isomorphism.\\
In particular, $\mathcal{U}'\otimes\mathcal{V}'\hookrightarrow(\mathcal{U}\otimes\mathcal{V})'$; $(\varphi\otimes\psi)(\mathbf{u}\otimes\mathbf{v})=\varphi(\mathbf{u})\psi(\mathbf{v})$ is an isomorphism when $\dim\mathcal{U},\dim\mathcal{V}<\infty$.
\begin{proof}
\begin{enumerate}[label=(\roman*)]
\setlength{\itemsep}{0pt}
\item By Theorem \ref{thm:tensor-product-of-linear-maps}, $(A,B)\mapsto A\otimes B$ is bilinear, hence induces $\Theta$.
\item Let $0\neq Z=\sum\limits_{i=1}^{r}A_i\otimes B_i$ with $r$ minimal. Then $A_i\neq0$ and $(B_i)$ is linearly independent.\\
Choose $\mathbf{u}$ with some $A_i\mathbf{u}\neq0$, after reindexing let $A_1\mathbf{u},\ldots,A_m\mathbf{u}$ be basis of $\operatorname{span}\{A_i\mathbf{u}:1\le i\le r\}$.\\
Absorb remaining coefficients into $B_i$ gives $C_1,\ldots,C_m$, not all zero, s.t. $\Theta(Z)(\mathbf{u}\otimes\mathbf{v})=\sum\limits_{i=1}^{m}(A_i\mathbf{u})\otimes C_i\mathbf{v}$.\\
If $\Theta(Z)=0$, Theorem \ref{thm:tensor-zero-test-and-decomposable-rank} gives $C_i\mathbf{v}=0$ for all $i,\mathbf{v}$, contradiction. Thus $\ker\Theta=0$.
\item If $\mathcal{U},\mathcal{V},\mathcal{X},\mathcal{Y}$ are finite-dimensional, then the domain and codomain of $\Theta$ have equal dimension by Theorem \ref{thm:tensor-bases-dimension-and-rank}, so injectivity implies bijectivity. Taking $\mathcal{X}=\mathcal{Y}=\mathbb{F}$, $\LL(\mathcal{U},\mathbb{F})=\mathcal{U}'$, $\LL(\mathcal{U}\otimes\mathcal{V},\mathbb{F})=(\mathcal{U}\otimes\mathcal{V})'$ gives the dual embedding and finite-dimensional isomorphism.
\end{enumerate}
\hlt{Structural Commentary}: Tensoring maps yields maps on tensor products.\\
\hlt{Key Insight}: The rule on decomposable tensors is enough, and finite dimensions make the embedding onto.
\end{proof}
\end{theorem}

\begin{proposition} \label{prop:scalar-extension-by-tensor-product} \hlt{Scalar Extension by Tensor Product}\\
Let $K/\mathbb{F}$ be a field extension and let $\mathcal{V}_{\mathbb{F}}$ be an $\mathbb{F}$-vector space.\\
Then $K\otimes_{\mathbb{F}}\mathcal{V}_{\mathbb{F}}$ is a $K$-vector space via $\alpha\cdot(\beta\otimes\mathbf{v})\defeq(\alpha\beta)\otimes\mathbf{v}$. With $\mu:\mathcal{V}_{\mathbb{F}}\to K\otimes_{\mathbb{F}}\mathcal{V}_{\mathbb{F}}$, $\mu(\mathbf{v})\defeq1\otimes\mathbf{v}$, every $\mathbb{F}$-linear $f:\mathcal{V}_{\mathbb{F}}\to\mathcal{Y}$ into a $K$-vector space $\mathcal{Y}$ admits a unique $K$-linear $\widetilde{f}:K\otimes_{\mathbb{F}}\mathcal{V}_{\mathbb{F}}\to\mathcal{Y}$ such that $\widetilde{f}\circ\mu=f$.
\begin{proof}
For $\alpha\in K$, $(\beta,\mathbf{v})\mapsto(\alpha\beta)\otimes\mathbf{v}$ is $\mathbb{F}$-bilinear, hence induces $M_\alpha\in\LL_{\mathbb{F}}(K\otimes_{\mathbb{F}}\mathcal{V}_{\mathbb{F}})$.\\
On decomposable tensors, $M_\alpha M_\lambda=M_{\alpha\lambda}$, $M_1=I$, and distributivity holds; since decomposable tensors span, these define a $K$-vector-space structure.\\
For $f$, set $g(\beta,\mathbf{v})\defeq\beta f(\mathbf{v})$. Then $g$ is $\mathbb{F}$-bilinear, so it induces $\widetilde{f}$ with $\widetilde{f}(\beta\otimes\mathbf{v})=\beta f(\mathbf{v})$. This is $K$-linear and gives $\widetilde{f}\circ\mu=f$. If $h$ is $K$-linear and $h\circ\mu=f$, then $h(\beta\otimes\mathbf{v})=\beta h(1\otimes\mathbf{v})=\beta f(\mathbf{v})$, so $h=\widetilde{f}$.\\
\hlt{Structural Commentary}: Scalar extension tensors in the larger field.\\
\hlt{Key Insight}: Tensoring with $K$ upgrades which scalars are allowed to act.
\end{proof}
\end{proposition}
